\chapter{Documentación de Visualización}
Este anexo presenta una descripción general de los \textit{scripts} que conforman el componente de visualización (la extensión de Blender). 
Las instrucciones de uso correspondientes se detallan en el Anexo~\ref{anx:manual_extraccion}.

\section{Requisitos}
Se requiere únicamente de Blender 5.2 LTS. Sin embargo, para modificar los scripts en editores
externos se puede utilizar un ambiente virtual con Python 3.13 con los siguientes
requisitos:

\begin{verbatim}
attrs==26.1.0
cattrs==26.1.0
certifi==2026.7.22
charset-normalizer==3.4.9
Cython==3.2.9
fake-bpy-module-5.2==20260730
idna==3.18
numpy==2.5.1
requests==2.34.2
typing_extensions==4.16.0
urllib3==2.7.0
zstandard==0.25.0
\end{verbatim}

Adicionalmente, se requiere el módulo de enlace para Python 3.13 de la biblioteca \textit{OpenVDB} (\texttt{pyopenvdb}). 
El procedimiento detallado de compilación e instalación se documenta en el Anexo~\ref{anx:compilacion_openvdb}.

\section{Organización del contenido}
El contenido del componente sigue una organización convencional para extensiones de Blender con la siguiente estructura
de directorio:
\begin{figure}[!htbp]
\dirtree{%
.1 TinyYoloV2Viewer/.
.2 assets/.
.3 library.blend.
.3 library.cats.txt.
.2 geometry\_nodes/.
.3 load.py.
.2 operators/.
.3 \_\_init\_\_.py.
.3 generate\_geometry.py.
.3 import\_json.py.
.2 properties/.
.3 \_\_init\_\_.py.
.3 graph\_properties.py.
.3 manifest\_properties.py.
.3 visualization\_properties.py.
.2 ui/.
.3 \_\_init\_\_.py.
.3 generate\_ui.py.
.2 utils/.
.3 utils.py.
.2 \_\_init\_\_.py.
.2 blender\_manifest.toml.
}
\caption{Estructura de archivos y directorios del paquete de la extensión \texttt{TinyYoloV2Viewer}.}
\label{fig:arbol_extension_blender}
\end{figure}

\section{Scripts}
La herramienta se distribuye como una extensión para Blender integrada por \textit{scripts} modulares en Python que exponen funcionalidad 
nativa en el editor, complementada por una biblioteca de recursos que contiene los elementos preconfigurados del sistema.

La estructura interna del paquete se organiza en los siguientes directorios:
\begin{itemize}
    \item \texttt{assets/}: Aloja el archivo contenedor \texttt{library.blend}, dentro del cual se definen los árboles de nodos de geometría 
    (\textit{Geometry Nodes}) y los materiales de sombreado requeridos por el complemento.
    \item \texttt{geometry\_nodes/}: Agrupa módulos auxiliares especializados en la vinculación, carga e instanciación procedimental de los 
    grupos de nodos contenidos en la biblioteca.
    \item \texttt{operators/}: Contiene los operadores de Blender (\texttt{bpy.types.Operator}) que ejecutan acciones sobre la escena, tales 
    como la lectura del manifiesto y la generación de geometrías volumétricas.
    \item \texttt{properties/}: Define las estructuras de datos (\texttt{bpy.types.PropertyGroup}) y propiedades personalizadas encargadas 
    de persistir el estado de la visualización y los metadatos de las capas importadas.
    \item \texttt{ui/}: Implementa los paneles gráficos (\texttt{bpy.types.Panel}) y menús desplegables integrados en la barra lateral de 
    la vista 3D.
    \item \texttt{utils/}: Provee funciones auxiliares para el registro de clases.
\end{itemize}

\subsection{Directorio \texttt{geometry\_nodes/}}
Este directorio contiene el módulo \texttt{load.py}, cuya finalidad consiste en gestionar la vinculación programática de los árboles de 
nodos de geometría procedentes de la biblioteca de recursos y parametrizar sus valores en tiempo de ejecución.

\subsubsection{\texttt{load.py}}
El módulo expone las siguientes funciones:

\begin{itemize}
    \item \texttt{load\_geonode\_group(group\_name)}: Importa un bloque de datos de tipo árbol de nodos (\texttt{NodeTree}) desde el archivo 
    \texttt{library.blend} a la escena activa. Recibe como argumento el identificador del grupo de nodos requerido.
    \item \texttt{deduplicate\_nodes()}: Identifica los bloques de nodos de geometría duplicados generados tras procesos de importación o 
    enlace redundantes (aquellos con sufijos automáticos del tipo \texttt{.001}, \texttt{.002}, etc.). Reasigna todas las referencias de 
    los modificadores al bloque de datos primario y purga las copias redundantes de la memoria.
    \item \texttt{deduplicate\_materials()}: Realiza una consolidación análoga a la deduplicación de nodos, purgando instancias redundantes 
    de materiales de sombreado y reasociando los enlaces al material original.
    \item \texttt{update\_input\_image(image\_path)}: Actualiza la textura de entrada asignada al material \texttt{YV2 Input Image}, vinculando 
    la imagen original procesada especificada en el archivo \texttt{manifest.json}. Recibe como parámetro la ruta absoluta o relativa de la imagen.
    \item \texttt{set\_node\_parameters(object, modifier\_name, parameters)}: Asigna valores a los campos de entrada de un modificador de nodos 
    de geometría (\textit{Geometry Nodes}) asociado a un objeto de la escena. Recibe como argumentos el objeto destino (\texttt{bpy.types.Object}), 
    el identificador del modificador (por defecto \texttt{"Geometry Nodes"}) y un diccionario de Python con las parejas clave-valor correspondientes 
    (\texttt{\{"Parámetro": valor\}}).
\end{itemize}

\subsection{Directorio \texttt{operators/}}
Este directorio agrupa los módulos encargados de registrar y ejecutar operadores nativos de Blender (\texttt{bpy.types.Operator}). Sus dos 
funcionalidades primarias consisten en la importación analítica del archivo \texttt{manifest.json} y la instanciación procedimental de las mallas 
tridimensionales y árboles de nodos a partir de dicha información.

\subsubsection{\texttt{import\_json.py}}
Este módulo implementa el operador \texttt{IMPORT\_OT\_load\_yolov2\_json}, responsable de la lectura y decodificación del archivo de manifiesto. 
Expone los siguientes métodos clave:

\begin{itemize}
    \item \texttt{execute(context)}: Lee el archivo \texttt{manifest.json} desde la ruta seleccionada en el explorador de archivos y persiste su 
    estructura estructurada en las propiedades personalizadas de la escena activa.
    \item \texttt{\_load\_properties()}: Parsea el contenido del manifiesto e inicializa las colecciones de propiedades de visualización asociadas 
    a cada tipología de capa.
    \item \texttt{menu\_func\_import(self, context)}: Función de integración que suscribe el operador al menú contextual global de importación de 
    Blender (\textit{File $\rightarrow$ Import $\rightarrow$ YOLOv2 Manifest (.json)}).
\end{itemize}

\subsubsection{\texttt{generate\_geometry.py}}
Este módulo implementa el operador \texttt{MESH\_OT\_create\_yolov2\_object}, encargado de construir la representación volumétrica y topológica 
de la red en la vista 3D. Define las siguientes rutinas de generación y configuración:

\begin{itemize}
    \item \texttt{execute(context)}: Orquesta la secuencia de construcción geométrica: crea los planos de visualización de entrada y salida, 
    reconstruye las capas intermedias, genera las aristas de conexión topológica y estructura los objetos resultantes en colecciones jerárquicas 
    dentro del \textit{Outliner}.
    \item \texttt{\_instance\_nodes(node\_group\_name, object\_name, context)}: Instancia una nueva entidad geométrica en la escena, le asigna un 
    modificador de nodos de geometría (\textit{Geometry Nodes}) y vincula el árbol de nodos especificado.
    \item \texttt{\_configure\_volume\_viewer(layer\_item)}: Parametriza el modificador \texttt{YV2 Volume Viewer} asociado al objeto de una capa 
    para vincular el archivo volumétrico \textit{OpenVDB} correspondiente y definir sus dimensiones nativas. Recibe como argumento una instancia 
    de \texttt{LayerItem}.
    \item \texttt{\_configure\_box\_viewer(target\_object, boxes\_path, parameters)}: Configura el modificador \texttt{YV2 Box Viewer} asignado 
    a la capa de salida para proyectar las cajas delimitadoras. Recibe como argumentos el objeto destino, la ruta de la cuadrícula \texttt{boxes.vdb} 
    y un diccionario con los umbrales de detección y filtrado.
    \item \texttt{\_configure\_input\_viewer(target\_object)}: Establece las dimensiones y texturas del modificador \texttt{YV2 Input Viewer} 
    sobre el objeto representativo de la imagen suministrada al modelo.
    \item \texttt{\_configure\_connections(connection\_object, parent\_layer, child\_layer)}: Parametriza el modificador \texttt{YV2 Volume Connection} 
    para calcular analíticamente las aristas de unión y el recorrido del \textit{kernel} convolucional entre dos instancias consecutivas de \texttt{LayerItem}.
    \item \texttt{\_layer\_parameters(layer\_item, manifest\_dir, layer\_data)}: Rellena las propiedades del contenedor \texttt{LayerItem} con los metadatos 
    estructurales procedentes del diccionario del manifiesto (resolución espacial, canal, paso de desplazamiento y rutas absolutas a los volúmenes).
    \item \texttt{\_get\_or\_create\_collection(name, parent\_collection)}: Evalúa la existencia de una colección en la escena y la instancia recursivamente 
    bajo la colección padre especificada si aún no ha sido registrada.
    \item \texttt{\_assign\_object\_to\_collection(object, target\_collection)}: Vincula la entidad geométrica a la colección destino y remueve sus referencias 
    de la colección raíz por defecto.
\end{itemize}

\subsection{Directorio \texttt{properties/}}
Este directorio contiene los módulos encargados de definir grupos de propiedades personalizadas (\texttt{bpy.types.PropertyGroup}) para garantizar la persistencia 
de los datos del modelo y la configuración de la vista dentro del archivo \texttt{.blend} entre distintas sesiones de trabajo. Las estructuras se dividen 
fundamentalmente en la persistencia del manifiesto deserializado y el control paramétrico de la visualización geométrica.

\subsubsection{\texttt{manifest\_properties.py}}
Define el contenedor de datos para almacenar el archivo de manifiesto a nivel de escena. Implementa la clase \texttt{ManifestDataSettings}, 
la cual expone dos propiedades de almacenamiento primarias:

\begin{itemize}
    \item \texttt{filepath}: Cadena de texto que registra la ruta del archivo \texttt{manifest.json} cargado.
    \item \texttt{raw\_data}: Cadena de texto que almacena la carga útil completa del JSON original en formato serializado, 
    solventando las restricciones de la API de Blender sobre tipos de datos anidados en las propiedades de escena.
\end{itemize}

Asimismo, la clase define los siguientes métodos de manipulación:
\begin{itemize}
    \item \texttt{get\_parsed\_json()}: Deserializa la cadena almacenada en \texttt{raw\_data} y retorna una estructura nativa de 
    diccionario en Python para su consumo por otros módulos.
    \item \texttt{set\_json(data\_string)}: Asigna y persiste el contenido crudo del manifiesto en las propiedades de la escena.
    \item \texttt{get\_manifest\_location()}: Retorna la ruta canónica del directorio contenedor del archivo de manifiesto para 
    resolver rutas relativas hacia las cuadrículas \textit{OpenVDB}.
\end{itemize}

\subsubsection{\texttt{visualization\_properties.py}}
Este módulo implementa las estructuras y rutinas de sincronización encargadas de controlar la apariencia visual de los volúmenes y 
la disposición espacial de la geometría.

En primer término, define la clase \texttt{LayerTypeProperty} (\texttt{PropertyGroup}) para agrupar los ajustes de renderizado 
aplicables a nivel de tipología de capa (e.g., \textit{Conv2d}, \textit{MaxPool2d}):
\begin{itemize}
    \item \texttt{name}: Cadena identificadora del tipo de capa.
    \item \texttt{color}: Vector de color RGBA asignado a la representación simplificada de la capa.
    \item \texttt{visibility}: Indicador booleano que determina si se fuerza la representación simplificados.
    \item \texttt{type\_visibility}: Indicador booleano que conmuta la visibilidad de la etiqueta tipográfica 
    que identifica la capa en la vista 3D.
\end{itemize}

Posteriormente, la clase \texttt{VisualizationDataProperties} centraliza los parámetros globales de disposición geométrica de la red:
\begin{itemize}
    \item \texttt{size\_scale}: Factor de escala aplicado a las dimensiones de altura y anchura (ejes $X$ e $Y$) en función de su resolución nativa.
    \item \texttt{layer\_depth}: Factor de escala longitudinal aplicado a la profundidad (eje $Z$).
    \item \texttt{layer\_gap}: Distancia de separación entre capas consecutivas.
    \item \texttt{layer\_max\_depth}: Valor máximo para la dimensión en el eje $Z$ de los volúmenes.
    \item \texttt{layer\_properties}: Colección de elementos de tipo \texttt{LayerTypeProperty} (\texttt{CollectionProperty}) 
    con las configuraciones particulares indexadas por tipo de capa.
\end{itemize}

Esta clase incorpora los siguientes métodos:
\begin{itemize}
    \item \texttt{create\_layer\_properties()}: Instancia e inicializa los registros de tipo \texttt{LayerTypeProperty} dentro 
    de la colección \texttt{layer\_properties}.
    \item \texttt{update\_all()}: Propaga los cambios de las propiedades globales a todos los objetos activos en la escena.
\end{itemize}

De forma complementaria, el módulo implementa funciones independientes de actualización reactiva (\textit{update callbacks}) que 
transfieren los valores modificados desde la interfaz hacia los nodos de geometría correspondientes:
\begin{itemize}
    \item \texttt{update\_colors(context)}: Propaga las modificaciones de color a los materiales de las cajas simplificadas.
    \item \texttt{update\_visibility(context)}: Actualiza el parámetro de visualización simplificada en los modificadores.
    \item \texttt{update\_dimensions(context)}: Recalcula y asigna los factores de escala espacial a los árboles de nodos de cada volumen.
    \item \texttt{update\_positions(context)}: Recalcula las coordenadas axiales y el desplazamiento longitudinal de las capas y conexiones 
    topológicas a lo largo del eje de propagación.
\end{itemize}

\subsubsection{\texttt{graph\_properties.py}}
Este módulo modela la estructura topológica y las propiedades geométricas de la red dentro del ecosistema de Blender. 
Implementa clases derivadas de \texttt{bpy.types.PropertyGroup} para representar nodos, conexiones y dependencias jerárquicas en 
forma de grafo acíclico dirigido (DAG), vinculando los metadatos de las capas con los objetos tridimensionales instanciados en la escena.

\paragraph{Clases auxiliares de colección}
Primero define clases auxiliares para el uso de colecciones de datos:
\begin{itemize}
    \item \texttt{StringItem}: Contenedor elemental que encapsula una cadena de texto en la propiedad 
    \texttt{value} (\texttt{StringProperty}), solventando la restricción de Blender que impide almacenar 
    listas nativas de cadenas dentro de una \texttt{CollectionProperty}. Se utiliza para listar identificadores de nodos padre e hijo.
    \item \texttt{ConnectionItem}: Contenedor para referencias a objetos de conexión. Expone la propiedad 
    de puntero \texttt{object} (\texttt{PointerProperty} de tipo \texttt{bpy.types.Object}), la cual almacena 
    la entidad de malla que contiene el modificador \texttt{YV2 Volume Connection}.
\end{itemize}

\paragraph{Clase \texttt{LayerItem}}
Representa individualmente una capa o nodo computacional de la red neuronal. Encapsula tanto los hiperparámetros 
operacionales del tensor como los enlaces a los datos volumétricos y a la entidad geométrica en la escena:
\begin{itemize}
    \item \texttt{name}: Identificador canónico y único de la capa dentro del grafo.
    \item \texttt{type}: Tipología de la operación computacional asociada (e.g., \textit{Conv2d}, \textit{MaxPool2d}, \textit{BatchNorm2d}).
    \item \texttt{dimensions}: Vector tridimensional (\texttt{FloatVectorProperty} de tamaño 3) que define las dimensiones métricas base ($X, Y, Z$) del prisma o volumen en la escena.
    \item \texttt{kernel\_size}: Tamaño espacial del núcleo de convolución o ventana de reducción ($K$).
    \item \texttt{stride}: Magnitud del paso de desplazamiento espacial aplicado por la operación.
    \item \texttt{ceiling}: Conmutador booleano que indica si se emplea la función techo (\textit{ceiling mode}) para el cálculo dimensional del tensor de salida.
    \item \texttt{resolution}: Resolución espacial bidimensional transversal de la capa ($H = W$).
    \item \texttt{slice\_count}: Cantidad total de planos longitudinales o canales ($C$) que componen el volumen.
    \item \texttt{file}: Ruta al archivo volumétrico \textit{OpenVDB} (\texttt{.vdb}) asociado a las activaciones de la capa.
    \item \texttt{object}: Puntero (\texttt{PointerProperty}) al objeto tridimensional de Blender (\texttt{bpy.types.Object}) asignado a la capa.
    \item \texttt{parents}: Colección (\texttt{CollectionProperty} de tipo \texttt{StringItem}) que lista los nombres de las capas emisoras o predecesoras directas.
    \item \texttt{children}: Colección (\texttt{CollectionProperty} de tipo \texttt{StringItem}) con los identificadores de las capas sucesoras o receptoras.
\end{itemize}

\paragraph{Clase \texttt{GraphProperty}}
Estructura central de la escena encargada de persistir el grafo computacional completo, sincronizar las dependencias 
jerárquicas y gestionar el ciclo de vida de los objetos tridimensionales generados.
\medskip

Expone las siguientes propiedades:
\begin{itemize}
    \item \texttt{layers}: Colección (\texttt{CollectionProperty} de tipo \texttt{LayerItem}) que contiene todos los nodos de la red indexados por nombre.
    \item \texttt{connections}: Colección (\texttt{CollectionProperty} de tipo \texttt{ConnectionItem}) que almacena los objetos de enlace e interpolación de 
    \textit{kernel} entre capas.
\end{itemize}

Asimismo, implementa los siguientes métodos de gestión y depuración:
\begin{itemize}
    \item \texttt{add\_connection(name, obj)}: Registra una nueva conexión topológica o actualiza una existente identificada por \texttt{name}, vinculando 
    el objeto de enlace \texttt{obj} (\texttt{bpy.types.Object}) y retornando el elemento instanciado.
    \item \texttt{add\_layer(name, type, obj)}: Da de alta un nuevo nodo en la colección \texttt{layers} especificando su nombre, tipo de módulo y puntero 
    al objeto en escena.
    \item \texttt{add\_node\_parent(parent, child)}: Establece una relación de precedencia dirigida insertando el identificador \texttt{parent} en la 
    colección \texttt{parents} del nodo \texttt{child}; si el nodo destino aún no existe en el registro, lo instancia dinámicamente.
    \item \texttt{set\_children\_nodes(graph\_node, children\_name)}: Actualiza de manera atómica las dependencias descendientes de una capa; vacía la 
    colección \texttt{children} del nodo \texttt{graph\_node}, registra bidireccionalmente cada elemento de la lista \texttt{children\_name} e invoca 
    internamente \texttt{add\_node\_parent}.
    \item \texttt{empty()}: Rutina de saneamiento que elimina la geometría generada en la escena para permitir su reconstrucción. Desvincula y destruye 
    los objetos de capas y conexiones de la base de datos (\texttt{bpy.data.objects}), purga los bloques de malla (\texttt{bpy.data.meshes}) que queden 
    sin usuarios huérfanos (\texttt{users == 0}) y vacía la colección \texttt{layers}.
\end{itemize}

\subsection{Directorio \texttt{ui/}}
Este directorio contiene el módulo \texttt{generate\_ui.py}, en el cual se define la clase de interfaz gráfica \texttt{VIEW3D\_PT\_Yolov2GeneratorPanel} 
(\texttt{bpy.types.Panel}). Su objetivo fundamental consiste en registrar e instanciar el panel de control interactivo dentro de la barra lateral 
(\textit{Sidebar} / región \textit{UI}) de la vista 3D (\textit{3D Viewport}), desplegando los controles para parametrizar las opciones de visualización, 
ajustar las dimensiones del grafo y accionar el botón de ejecución del operador de generación geométrica (\texttt{MESH\_OT\_create\_yolov2\_object}).

\section{Recursos}
El archivo \texttt{library.blend} contiene una librería de recursos compuesto principalmente por los grupos de nodos de geometría y materiales utilizados
por la extensión. A continuación se describen estos.

\subsection{Nodos de Geometría}
Los grupos de nodos principales utilizados para la visualización son \texttt{YV2 Volume Viewer}, \texttt{YV2 Box Viewer}, \texttt{YV2 Input Viewer}
y \texttt{YV2 Volume Connection}, los cuales pueden ser instanciados como modificadores de geometría. Otros grupos de nodos son incluidos, aunque
su uso es meramente auxiliar.

\subsubsection{YV2 Volume Viewer}
Constituye el grupo de nodos primario del sistema de visualización. Su objetivo consiste en reconstruir los archivos de cuadrícula \textit{OpenVDB} (\texttt{.vdb}) 
mediante una secuencia de planos coplanares que representan las rebanadas (\textit{slices}) del volumen. El nodo expone los siguientes parámetros de configuración:

\begin{itemize}
    \item \texttt{Volume File}: Ruta absoluta o relativa del archivo volumétrico \textit{OpenVDB} a cargar.
    \item \texttt{Type}: Cadena de texto descriptiva que indica el tipo de capa o metadatos informativos a desplegar.
    \item \texttt{View Type}: Conmutador booleano (\textit{toggle}) para habilitar o suspender la visibilidad de la etiqueta de texto definida en \texttt{Type}.
    \item \texttt{Simplified Volume}: Conmutador que sustituye el conjunto de planos de rebanadas por un prisma envolvente sólido. Se recomienda 
    mantener activa esta vista simplificada en aquellas capas que no se encuentren bajo inspección activa a fin de reducir la carga de cómputo geométrico en la vista 3D.
    \item \texttt{Color}: Valor cromático asignado a la malla simplificada de la capa.
    \item \texttt{World Size}: Dimensiones globales del volumen en unidades métricas de Blender (\textit{Blender Units}). Es imperativo modificar 
    la escala del volumen mediante este parámetro en lugar de recurrir a transformaciones afines de objeto (\textit{Object Transform}), preservando 
    así la correctitud espacial de la visualización de las conexiones.
    \item \texttt{Resolution}: Resolución bidimensional de las rebanadas ($H \times W$), asumiendo una sección transversal cuadrada.
    \item \texttt{Slice Count}: Cantidad total de rebanadas contenidas a lo largo del eje de profundidad longitudinal ($C$).
    \item \texttt{Slice Selection}: Índice ordinal de la rebanada que se desea aislar; atenúa u oculta la geometría de las rebanadas anteriores en la secuencia.
    \item \texttt{Single Slice}: Conmutador que restringe la visualización exclusivamente a la rebanada indicada en \texttt{Slice Selection}, descartando la 
    renderización de las rebanadas anteriores y posteriores para optimizar el rendimiento y permitir la inspección transversal aislada.
\end{itemize}

Este grupo de nodos instancia una cuadrícula de polígonos (planos bidimensionales subdivididos) por cada canal del tensor volumétrico,
operando cada uno como una rebanada transversal del mismo. Cada cuadrícula se genera con una densidad de subdivisiones equivalente a la 
resolución nativa del volumen a representar, de modo que cada cara poligonal corresponde biunívocamente con un vóxel en la cuadrícula de datos discretos.
A partir de las coordenadas del centroide de cada cara, se muestrea el campo volumétrico continuo del archivo \textit{OpenVDB} subyacente y se almacena 
la intensidad resultante como un atributo de dominio de cara (\textit{face domain attribute}), permitiendo que el shader mapee dicho escalar a una rampa 
cromática. De forma complementaria, el árbol de nodos evalúa la visibilidad condicional de las cuadrículas para conmutar o descartar rebanadas en función 
de los parámetros de configuración activos.
\medskip

Asimismo, este grupo de nodos escribe atributos geométricos persistentes en la malla para facilitar el intercambio analítico de datos y la 
sincronización con otros grupos de nodos (como \texttt{YV2 Volume Connection}):

\begin{itemize}
    \item \texttt{grid\_color}: Atributo escalar asignado por cara en cada rebanada. Almacena el valor muestreado del volumen en la 
    coordenada discreta asociada, consumido directamente por el material de sombreado.
    \item \texttt{tensor\_size}: Vector tridimensional con las dimensiones métricas globales del volumen en la escena 
    (equivalente al parámetro de entrada \texttt{World Size}).
    \item \texttt{tensor\_resolution}: Resolución espacial bidimensional transversal ($H = W$) del tensor (equivalente a \texttt{Resolution}).
    \item \texttt{tensor\_layer\_count}: Cantidad total de planos longitudinales o canales ($C$) que componen el volumen (equivalente a \texttt{Slice Count}).
    \item \texttt{tensor\_layer\_selection}: Índice entero del canal o rebanada activa seleccionada para inspección individual (equivalente a \texttt{Slice Selection}).
\end{itemize}

\subsubsection{YV2 Box Viewer}
Este grupo de nodos de geometría se utiliza para visualizar las cajas delimitadoras (\textit{bounding boxes}) generadas por la red. Su función principal 
consiste en decodificar los atributos de detección contenidos en la cuadrícula volumétrica de \textit{OpenVDB} (\texttt{boxes.vdb}) y reconstruir proceduralmente 
la geometría de las cajas (posiciones espaciales, dimensiones normalizadas y etiquetas de clase) proyectándolas sobre la imagen de entrada del modelo.
\medskip

El nodo expone los siguientes parámetros:
\begin{itemize}
    \item \texttt{Box File}: Ruta al archivo \textit{OpenVDB} (\texttt{boxes.vdb}) que contiene los datos matriciales codificados de las predicciones.
    \item \texttt{World Size}: Dimensiones globales de la visualización en unidades métricas de Blender. Únicamente los componentes de los ejes $X$ e $Y$ 
    (anchura y altura, respectivamente) modifican la escala proyectiva del plano de visualización.
    \item \texttt{Class Labels}: Cadena de texto que contiene el listado de las etiquetas categóricas de la red, separadas por comas.
    \item \texttt{Box Count}: Número total de cajas delimitadoras generadas por la red antes o después del filtrado.
    \item \texttt{Box Selection}: Selector numérico que controla el régimen de despliegue de las predicciones:
    \begin{itemize}
        \item \texttt{-1}: Visualiza todas las cajas candidatas que superen los umbrales configurados en \texttt{Confidence} y \texttt{Class Probability}.
        \item \texttt{0} a \texttt{Box Count - 1}: Aísla e ilustra una única caja de predicción conforme a su índice ordinal.
        \item \texttt{Box Count}: Fuerza la visualización de la totalidad de las cajas generadas sin aplicar filtrado.
    \end{itemize}
    \item \texttt{Confidence}: Umbral de confianza de detección (\textit{objectness score}). Filtra las cajas cuya certeza calculada sea inferior a este valor.
    \item \texttt{Class Probability}: Umbral de probabilidad condicional de clase. Descarta las predicciones cuya confianza categórica no alcance dicha magnitud.
\end{itemize}
El grupo de nodos instancia un plano que proyecta la imagen de entrada junto con una malla alámbrica cuadrangular por cada caja delimitadora candidata. 
A partir del índice de cada instancia y del esquema de direccionamiento discreto codificado previamente en la cuadrícula volumétrica de 
\textit{OpenVDB} (\texttt{boxes.vdb}), el árbol recupera el vector de características de cada predicción.

Los índices correspondientes a este vector estructurado son:
\begin{itemize}
    \item $0$: Coordenada central horizontal ($x$).
    \item $1$: Coordenada central vertical ($y$).
    \item $2$: Anchura de la caja delimitadora ($w$).
    \item $3$: Altura de la caja delimitadora ($h$).
    \item $4$: Puntuación de objetualidad (\textit{objectness score}).
    \item $5$: Probabilidad condicional de la clase predicha ($P(\text{clase})$).
    \item $6$: Identificador categórico de la clase (\textit{class ID}).
\end{itemize}
A partir de estas componentes métricas normalizadas, se calculan la translación espacial y la escala local del plano alámbrico de forma relativa al 
plano de referencia de la imagen. Adicionalmente, el árbol genera geometrías de texto tridimensional a partir de las cadenas suministradas en 
el parámetro \texttt{Class Labels} e instancia la etiqueta textual sobre el borde superior de cada caja en función de su identificador de clase.
\medskip 

Asimismo, el identificador entero de clase se normaliza al intervalo continuo $[0, 1]$ respecto al total de categorías disponibles en el conjunto 
de datos; este valor escalar se transfiere como atributo al shader de la malla para modular una rampa de color que diferencia visualmente cada categoría 
inferida. Lógica booleana complementaria evalúa los umbrales de confianza y filtros de visibilidad para descartar o mostrar las cajas en tiempo real.
\medskip

Este grupo de nodos escribe los siguientes atributos geométricos en la malla para su consumo por el material y la comunicación con otros modificadores:
\begin{itemize}
    \item \texttt{class}: Identificador de clase normalizado en el intervalo $[0, 1]$, consumido por el material para asignar una 
    signatura cromática unívoca a la caja y a su etiqueta textual.
    \item \texttt{tensor\_size}: Vector tridimensional con las dimensiones métricas globales del plano de visualización en el espacio de Blender 
    (equivalente al parámetro de entrada \texttt{World Size}).
\end{itemize}

\subsubsection{YV2 Input Viewer}
Este grupo de nodos se utiliza para proyectar e inspeccionar la imagen de entrada suministrada a la red. Expone los siguientes parámetros:

\begin{itemize}
    \item \texttt{World Size}: Dimensiones espaciales del plano de visualización en unidades métricas de Blender (\textit{Blender Units}).
    \item \texttt{Resolution}: Resolución bidimensional de la imagen de entrada ($H \times W$), asumiendo una relación de aspecto cuadrada.
    \item \texttt{View Channels}: Conmutador booleano que determina si se descomponen y despliegan los canales individuales RGB 
    (rojo, verde y azul) en planos separados, en adición a la composición original a color.
\end{itemize}
Este grupo de nodos instancia cuatro planos bidimensionales que proyectan la imagen de entrada, donde cada plano aplica un factor cromático 
multiplicativo en el sombreador (\textit{shader}) para descomponer e ilustrar sus canales correspondientes.
\medskip

Asimismo, este grupo de nodos escribe atributos geométricos en la malla para facilitar la transferencia de datos e interoperabilidad con 
otros árboles de nodos de la escena:

\begin{itemize}
    \item \texttt{tensor\_size}: Vector tridimensional con las dimensiones métricas del plano en el espacio de trabajo 
    (equivalente al parámetro de entrada \texttt{World Size}).
    \item \texttt{tensor\_resolution}: Resolución espacial bidimensional transversal ($H = W$) asignada a la imagen 
    (equivalente a \texttt{Resolution}).
    \item \texttt{tensor\_layer\_count}: Cantidad total de planos o canales longitudinales instanciados; se fija con 
    un valor constante de $4$.
    \item \texttt{tensor\_layer\_selection}: Índice de la rebanada o canal seleccionado para visualización activa; se 
    inicializa con un valor de $0$.
\end{itemize}

\subsubsection{YV2 Volume Connection}
Este grupo de nodos de geometría se encarga de representar visualmente las conexiones topológicas entre las distintas capas de la red, así como las operaciones 
matemáticas de convolución o reducción espacial que ocurren entre ellas. 
\medskip

El modificador expone los siguientes parámetros:
\begin{itemize}
    \item \texttt{Input Layer}: Objeto tridimensional que representa la capa previa o nodo emisor en el grafo computacional. Requiere que dicho objeto contenga una 
    instancia activa del modificador \texttt{YV2 Volume Viewer} debidamente parametrizada.
    \item \texttt{Output Layer}: Objeto tridimensional que representa la capa sucesora o nodo receptor de la relación. De igual forma, asume la presencia y correcta 
    configuración del modificador \texttt{YV2 Volume Viewer}.
    \item \texttt{View Convolution}: Conmutador booleano que habilita la renderización de la geometría correspondiente al \textit{kernel} convolucional y sus líneas de proyección.
    \item \texttt{T}: Parámetro normalizado de interpolación temporal para la animación interactiva de la convolución ($t \in [0, 1]$). Controla el barrido secuencial del 
    \textit{kernel} sobre el volumen de entrada y el punto resultante en el plano de salida, iniciando en la esquina superior izquierda ($0.0$) y concluyendo en la esquina inferior derecha ($1.0$).
    \item \texttt{Stride}: Paso de desplazamiento espacial (\textit{stride}) del núcleo en los ejes de altura y anchura. Omite posiciones intermedias de la entrada 
    según el incremento configurado. No realiza comprobaciones de consistencia dimensional entre tensores.
    \item \texttt{Ceiling Mode}: Conmutador booleano que selecciona el criterio de redondeo al calcular las dimensiones de salida si el paso fracciona el recorrido 
    sobre las filas y columnas. Aplica la función suelo (\textit{floor}) cuando está inactivo y la función techo (\textit{ceiling}) cuando está habilitado 
    (análogo al parámetro homónimo en operaciones como \textit{MaxPool2d}).
    \item \texttt{Operation Mode}: Selector categórico del régimen de mapeo de canales entre volúmenes:
    \begin{itemize}
        \item \textit{One to One} (canal a canal): Procesa cada canal de entrada de forma independiente sobre su canal análogo de salida (e.g., submuestreo por \textit{MaxPool2d}).
        \item \textit{Many to One} (volumen a plano/canal): Modela la agregación de todos los canales de entrada para proyectar una única rebanada o 
        mapa de activación resultante (convolución estándar).
        \item \textit{Many to Many}: Representación generalizada del flujo volumétrico completo entre el bloque entrante y saliente.
    \end{itemize}
\end{itemize}
A partir de la resolución (\texttt{tensor\_resolution}) y las dimensiones espaciales (\texttt{tensor\_size}) de los volúmenes enlazados, junto con los parámetros 
\texttt{Stride} y \texttt{Ceiling Mode}, este grupo de nodos calcula analíticamente las coordenadas tridimensionales de cada vóxel en las representaciones 
volumétricas conectadas. La trayectoria espacial toma como origen el vóxel superior izquierdo de la rebanada seleccionada y como destino final el vóxel inferior 
derecho. Mediante el parámetro de interpolación normalizado \texttt{T}, se interpola linealmente a lo largo de este barrido matricial, recorriendo las celdas 
de izquierda a derecha y de arriba hacia abajo. Apoyándose en los atributos \texttt{tensor\_layer\_count} y \texttt{tensor\_layer\_selection}, determina 
adicionalmente la cota longitudinal (profundidad en el eje $Z$) del canal bajo inspección activa. Cabe destacar que el modificador asume la correspondencia 
dimensional entre tensores y no realiza validaciones internas para verificar que las dimensiones de entrada y salida concuerden con los hiperparámetros de la capa conectada.
\medskip

Con base en la posición discreta del vóxel evaluado, el árbol genera dos prismas alámbricos que delimitan espacialmente el núcleo convolucional 
(\textit{kernel}) y su vecindad receptiva, dimensionándolos y orientándolos de acuerdo con el modo de convolución configurado en \texttt{Operation Mode}. 
De forma complementaria, proyecta una réplica del contorno del núcleo directamente sobre las rebanadas actualmente seleccionadas para delimitar visualmente 
los valores tensoriales sometidos a la operación en ese instante.
\medskip

Asimismo, el grupo de nodos genera aristas de interconexión que unen las caras perimetrales homólogas de ambos volúmenes.

\subsubsection{Nodos auxiliares}
A continuación, se describen los grupos de nodos auxiliares. Estos bloques no están diseñados para instanciarse directamente como modificadores independientes, 
sino que actúan como subrutinas modulares compartidas por los grupos principales de \textit{Geometry Nodes}, facilitando la reutilización de lógica matemática 
y la encapsulación funcional:

\begin{itemize}
    \item \texttt{YV2 Ceiling Mode}: A partir de un conmutador booleano y un escalar continuo, aplica una función de redondeo hacia abajo 
    (\textit{floor}, si es \texttt{False}) o hacia arriba (\textit{ceil}, si es \texttt{True}).
    \item \texttt{YV2 Class Name Mesh}: Genera mallas poligonales de texto tridimensional para cada categoría a partir de la lista de 
    etiquetas de clase, indexándolas conforme a sus respectivos identificadores de clase.
    \item \texttt{YV2 Convolution Depth}: Determina la profundidad (eje $Z$) correspondiente a una rebanada inspeccionada, en función 
    del modo de conexión seleccionado, la capa activa y el canal asignado.
    \item \texttt{YV2 Create Line}: Genera una primitiva geométrica lineal entre dos posiciones.
    \item \texttt{YV2 Even Kernel Correction}: Ajusta el centroide y desfase espacial de la vecindad del núcleo convolucional cuando la 
    dimensión del \textit{kernel} corresponde a un valor par.
    \item \texttt{YV2 Get Box Info}: Decodifica matricialmente el vector de características de una predicción a partir de la cuadrícula 
    discreta almacenada en el archivo \texttt{boxes.vdb}.
    \item \texttt{YV2 Get Layer Info}: Extrae los atributos geométricos persistentes de un objeto (\texttt{tensor\_size}, 
    \texttt{tensor\_resolution}, \texttt{tensor\_layer\_count} y \texttt{tensor\_layer\_selection}), 
    así como su matriz de transformación espacial.
    \item \texttt{YV2 Get Volume Corner}: Calcula los vértices extremos de la caja envolvente de un volumen a partir de sus dimensiones 
    nativas y su matriz de transformación. Dispone de conmutadores booleanos para seleccionar vértices específicos, retornando por defecto 
    la esquina frontal superior izquierda.
    \item \texttt{YV2 Iterate 2D Grid}: Computa las coordenadas espaciales requeridas para recorrer discretamente la retícula de una rebanada, 
    calculando la posición euclidiana de una celda en función del parámetro normalizado \texttt{T}, la resolución, las dimensiones métricas, 
    el valor de \textit{stride} y el modo de redondeo.
    \item \texttt{YV2 Kernel}: Construye la representación alámbrica del núcleo volumétrico y de su proyección sobre la rebanada activa evaluada.
    \item \texttt{YV2 Layer Box Connection}: Genera las aristas y polígonos de conexión entre las caras opuestas de dos volúmenes adyacentes.
    \item \texttt{YV2 Volume Corner Offset}: Determina el vector de desplazamiento hacia el centroide de la celda de origen 
    (esquina superior izquierda) en la rebanada de un volumen.
    \item \texttt{YV2 Wireframe Cube}: Genera una estructura alámbrica rectangular parametrizable para delimitar prismas cúbicos 
    o cuadriláteros bidimensionales.
\end{itemize}

\subsection{Materiales de sombreado}
A continuación, se describen los árboles de materiales (\textit{shaders}) empleados por la extensión para la renderización 
volumétrica y geométrica de la red en la escena. Todos los materiales implementan un modelo de sombreado predominantemente emisivo 
(\textit{Emission}) para asegurar una lectura homogénea y nítida de los datos independientes de fuentes de iluminación externa:

\begin{itemize}
    \item \texttt{YV2 Volume Slice}: Asignado a los planos bidimensionales que componen las rebanadas del tensor. Muestrea el atributo de cara 
    \texttt{grid\_color} proveniente de la cuadrícula de \textit{OpenVDB} y lo proyecta sobre una rampa cromática para generar una 
    representación en mapa de calor de las activaciones.
    \item \texttt{YV2 Simplified Volume}: Utilizado en los prismas poligonales durante el modo de visualización simplificado. Lee el vector 
    de color \texttt{volume\_color} asignado por tipo de capa e interpola sus coordenadas UV locales para proyectar un gradiente suave en las caras de la envolvente.
    \item \texttt{YV2 Input Image}: Asignado a los planos de la imagen de entrada suministrada a la red. Lee el atributo escalar o vectorial 
    \texttt{channel\_color} para aislar y modular el matiz cromático de los canales individuales proyectados.
    \item \texttt{YV2 Box Class}: Empleado para la diferenciación cromática de las cajas delimitadoras y sus etiquetas de texto 
    correspondientes. Muestrea el atributo normalizado \texttt{class} (comprendido en el intervalo $[0, 1]$) y modula directamente 
    el canal de matiz (\textit{hue}) en el espacio HSV, asignando una signatura cromática discreta y determinista a cada categoría del modelo.
    \item \texttt{YV2 Volume BB}: Material emisivo acromático (blanco puro) utilizado de forma transversal para representar las aristas 
    de conexión topológica, los prismas alámbricos de las cajas envolventes y las mallas tipográficas que identifican a cada capa.
\end{itemize}

\section{Pruebas de rendimiento}
En adición a la extensión se proveen scripts para la evaluación del rendimiento del mismo. Este esta compuesto de dos scripts de Python:
\texttt{run.py} y \texttt{benchmark.py}.

\subsection{\texttt{run.py}}
Ejecuta una instancia de Blender en modo desatendido (\textit{headless}). Esta rutina cuantifica de forma aislada 
los tiempos de ejecución correspondientes a la importación del manifiesto, la instanciación procedimental de la 
geometría y el renderizado de un fotograma de la escena.
\medskip

Este script contiene las siguientes funciones:
\begin{itemize}
    \item \texttt{run(input\_directory)}: Ejecuta el flujo integral de importación de datos, generación procedimental de geometría y 
    renderizado secuencial a partir de los datos contenidos en el directorio especificado.
    \item \texttt{warmup()}: Realiza una pasada preliminar de renderizado (\textit{warm-up}) para compilar los sombreadores en la GPU 
    y precargar las estructuras de datos y cuadrículas volumétricas en memoria caché.
    \item \texttt{run\_render()}: Ejecuta el renderizado del fotograma actual en la escena activa.
    \item \texttt{set\_camera()}: Ajusta la posición, orientación y encuadre del objeto de cámara activa para abarcar la totalidad de 
    la geometría del grafo generado dentro del campo visual (\textit{frustum}).
    \item \texttt{save\_render\_result(output\_path, filename)}: Guarda la imagen resultante del búfer de renderizado en la ruta de 
    destino y con el nombre de archivo indicados.
\end{itemize}

Adicionalmente, se expone la variable global \texttt{SAVE\_RENDER} para determinar si el proceso debe persistir las imágenes generadas 
en disco. Al concluir, el \textit{script} genera en el directorio de entrada un archivo JSON denominado 
\texttt{blender\_headless\_benchmark.json} con la telemetría individual de la ejecución, registrando las siguientes métricas:

\begin{itemize}
    \item \texttt{Import Runtime}: Tiempo transcurrido durante la carga y decodificación del archivo de manifiesto (en segundos).
    \item \texttt{Geometry Generation Runtime}: Tiempo requerido para la instanciación de mallas procedimentales y 
    la vinculación de cuadrículas volumétricas (en segundos).
    \item \texttt{Simple Render Runtime}: Tiempo de cómputo para el renderizado del fotograma configurado en modo de 
    prismas envolventes simplificados (en segundos).
    \item \texttt{One Layer Render Runtime}: Tiempo de renderizado habilitando la visualización detallada de las capas 
    convolucionales (\textit{Conv2d}, en segundos).
    \item \texttt{Two Layer Render Runtime}: Tiempo de renderizado habilitando la visualización detallada adicional de 
    las capas de reducción espacial (\textit{MaxPool2d}, en segundos).
    \item \texttt{Three Layer Render Runtime}: Tiempo de renderizado habilitando la visualización detallada adicional 
    de las capas de normalización por lotes (\textit{BatchNorm2d}, en segundos).
    \item \texttt{Complex Render Runtime}: Tiempo de renderizado de la red completa con la visualización volumétrica 
    detallada de todas las capas activas (en segundos).
    \item \texttt{Render Engine}: Identificador del motor de renderizado empleado durante la prueba (e.g., \textit{Eevee} o \textit{Cycles}).
    \item \texttt{Input Directory}: Ruta del directorio de origen de los datos evaluados.
\end{itemize}

\subsection{\texttt{benchmark.py}}
Este módulo se encarga de orquestar la ejecución por lotes (\textit{batch processing}) de las pruebas de rendimiento y estabilidad en Blender, consolidando las métricas de telemetría resultantes en un único archivo de balance. 
\medskip

La rutina invoca iterativamente al \textit{script} de automatización sin interfaz gráfica (\textit{headless}) \texttt{run.py} sobre cada uno de los subdirectorios de datos extraídos ubicados en la ruta definida por la constante global \texttt{INPUT\_PATH}. Posteriormente, recopila los registros generados por cada ejecución individual y los agrega en un archivo estructurado dentro del directorio especificado, calculando métricas estadísticas globales para el lote procesado:

\begin{itemize}
    \item \texttt{Average Import Runtime}: Media aritmética del tiempo de lectura, carga y decodificación del 
    archivo \texttt{manifest.json} (en segundos).
    \item \texttt{Average Geometry Generation Runtime}: Media aritmética del tiempo empleado por el operador en 
    instanciar la geometría procedimental y vincular los árboles de nodos (en segundos).
    \item \texttt{Average Simple Render Runtime}: Tiempo promedio de renderizado por cuadro empleando la representación 
    volumétrica simplificada mediante prismas envolventes (en segundos).
    \item \texttt{Average One Layer Render Runtime}: Tiempo promedio de renderizado habilitando la inspección 
    detallada de un único tipo de capa convolucional (\textit{Conv2d}, en segundos).
    \item \texttt{Average Two Layer Render Runtime}: Tiempo promedio de renderizado habilitando la inspección 
    volumétrica detallada de dos tipos de capa (\textit{Conv2d} y \textit{MaxPool2d}, en segundos).
    \item \texttt{Average Three Layer Render Runtime}: Tiempo promedio de renderizado habilitando la inspección 
    detallada de tres tipos de capa (\textit{Conv2d}, \textit{MaxPool2d} y \textit{BatchNorm2d}, en segundos).
    \item \texttt{Average Complex Render Runtime}: Tiempo promedio de renderizado en modo de visualización completa, 
    con la totalidad de capas intermedias y rebanadas volumétricas activas (en segundos).
    \item \texttt{Total Items}: Cantidad neta de muestras o ejecuciones individuales procesadas a lo largo de la corrida.
    \item \texttt{Runs}: Estructura jerárquica (diccionario) que almacena la telemetría detallada de cada corrida 
    individual, indexada por el identificador o ruta del subdirectorio de datos correspondiente.
\end{itemize}
 %% AGREGAR BENCHMARK